% GENERATED FILE. Edit canonical lessons, then run npm run generate:books.
% canonical-source-hash: fnv1a64:2710d35dec2a33c7
% canonical-lessons: BN-C384-rel, BN-C384-kamra, BN-C384-barth, BN-C384-lokal, BN-C384-huisol-deoya

\chapter{Rails and Carriages}
\label{ch:bn-rails-and-carriages}
\hlchaptermodality{384}

\begin{chapteropening}
\textbf{By the end of this chapter:} \emph{I can talk about the railway, a compartment and a berth, and the whistle as a train leaves.}

Complete the last lesson of chapter 384: I can talk about the railway, a compartment and a berth, and the whistle as a train leaves.
\end{chapteropening}

\section[\texorpdfstring{\bn{রেল} (rel)}{রেল (rel)}]{\bn{রেল} (rel) — a railway; the railways}

\label{lesson:BN-C384-rel}



\begin{quote}
Before the new one: say the Bengali for a female singer, then the Bengali for a female dancer.
\end{quote}

\subsection*{You'll want to know: \bn{রেল}}
\textbf{\bn{রেল}} — \emph{rel} — \textquotedblleft{}a railway; the railways\textquotedblright{}.

\begin{quote}
\bn{আমার} \bn{বাবা} \bn{রেলে} \bn{কাজ} \bn{করতেন।}
\end{quote}

\begin{quote}
\emph{āmār bābā rele kāj kôrten.}
\end{quote}

\begin{quote}
My father used to work on the railways.
\end{quote}

\bn{রেলে}, on the railways; \bn{রেলের}, of the railway: \bn{রেলের} \bn{টিকিট}, a railway ticket.

\begin{cousinweb}
\textbf{\bn{রেল}} (\emph{rel}) is English \emph{rail}, taken over in the age of steam for the line and then for the whole railway.
\end{cousinweb}

\begin{grammarlens}[title={What you've built}]
One more word for getting around.
\end{grammarlens}

\subsection*{Your turn}
\begin{itemize}
\raggedright
  \item \emph{Say it:} \emph{rel}
  \item \emph{Say it:} \emph{rel}, once more
  \item \emph{Say it:} \emph{āmār bābā rele kāj kôrten.} — My father used to work on the railways.
  \item \emph{Recall:} say the Bengali for a female singer, then the Bengali for a female dancer, then say \emph{rel} again
\end{itemize}

\begin{tcolorbox}[breakable,colback=teal!4,colframe=teal!35!black,title={Before you move on}]
What does \bn{রেল} mean? (\textquotedblleft{}A railway; the railways\textquotedblright{}.) Say it once more.
\end{tcolorbox}
\section[\texorpdfstring{\bn{কামরা} (kāmrā)}{কামরা (kāmrā)}]{\bn{কামরা} (kāmrā) — a compartment (of a train); a chamber}

\label{lesson:BN-C384-kamra}



\begin{quote}
Before the new one: say the Bengali for a female dancer, then the Bengali for a railway.
\end{quote}

\subsection*{You'll want to know: \bn{কামরা}}
\textbf{\bn{কামরা}} — \emph{kāmrā} — \textquotedblleft{}a compartment (of a train); a chamber\textquotedblright{}.

\begin{quote}
\bn{ট্রেনের} \bn{এই} \bn{কামরাটা} \bn{খালি।}
\end{quote}

\begin{quote}
\emph{ṭrener ei kāmrāṭā khāli.}
\end{quote}

\begin{quote}
This compartment of the train is empty.
\end{quote}

With the classifier: \bn{কামরাটা}, the compartment; \bn{মহিলাদের} \bn{কামরা}, the ladies' compartment.

\begin{cousinweb}
\textbf{\bn{কামরা}} (\emph{kāmrā}) is usually traced to Portuguese \emph{câmara}, \textquotedblleft{}a room, a chamber\textquotedblright{}, one of the words the Portuguese left along the Bengal coast; in Bengali it settled on the rooms of a train.
\end{cousinweb}

\begin{grammarlens}[title={What you've built}]
One more word for getting around.
\end{grammarlens}

\subsection*{Your turn}
\begin{itemize}
\raggedright
  \item \emph{Say it:} \emph{kāmrā}
  \item \emph{Say it:} \emph{kāmrā}, once more
  \item \emph{Say it:} \emph{ṭrener ei kāmrāṭā khāli.} — This compartment of the train is empty.
  \item \emph{Recall:} say the Bengali for a female dancer, then the Bengali for a railway, then say \emph{kāmrā} again
\end{itemize}

\begin{tcolorbox}[breakable,colback=teal!4,colframe=teal!35!black,title={Before you move on}]
What does \bn{কামরা} mean? (\textquotedblleft{}A compartment (of a train); a chamber\textquotedblright{}.) Say it once more.
\end{tcolorbox}
\section[\texorpdfstring{\bn{বার্থ} (bārth)}{বার্থ (bārth)}]{\bn{বার্থ} (bārth) — a berth (on a train)}

\label{lesson:BN-C384-barth}



\begin{quote}
Before the new one: say the Bengali for a railway, then the Bengali for a compartment.
\end{quote}

\subsection*{You'll want to know: \bn{বার্থ}}
\textbf{\bn{বার্থ}} — \emph{bārth} — \textquotedblleft{}a berth (on a train)\textquotedblright{}.

\begin{quote}
\bn{রাতের} \bn{ট্রেনে} \bn{আমার} \bn{বার্থটা} \bn{ওপরে} \bn{ছিল।}
\end{quote}

\begin{quote}
\emph{rāter ṭrene āmār bārthṭā opôre chhilô.}
\end{quote}

\begin{quote}
On the night train my berth was the upper one.
\end{quote}

\begin{cousinweb}
\textbf{\bn{বার্থ}} (\emph{bārth}) is borrowed from English \emph{berth}, a sleeping place on a ship or train; the English th is said with \bn{থ}.
\end{cousinweb}

\begin{grammarlens}[title={What you've built}]
One more word for getting around.
\end{grammarlens}

\subsection*{Your turn}
\begin{itemize}
\raggedright
  \item \emph{Say it:} \emph{bārth}
  \item \emph{Say it:} \emph{bārth}, once more
  \item \emph{Say it:} \emph{rāter ṭrene āmār bārthṭā opôre chhilô.} — On the night train my berth was the upper one.
  \item \emph{Recall:} say the Bengali for a railway, then the Bengali for a compartment, then say \emph{bārth} again
\end{itemize}

\begin{tcolorbox}[breakable,colback=teal!4,colframe=teal!35!black,title={Before you move on}]
What does \bn{বার্থ} mean? (\textquotedblleft{}A berth (on a train)\textquotedblright{}.) Say it once more.
\end{tcolorbox}
\section[\texorpdfstring{\bn{লোকাল} (lokāl)}{লোকাল (lokāl)}]{\bn{লোকাল} (lokāl) — a local train; a stopping train}

\label{lesson:BN-C384-lokal}



\begin{quote}
Before the new one: say the Bengali for a compartment, then the Bengali for a berth.
\end{quote}

\subsection*{You'll want to know: \bn{লোকাল}}
\textbf{\bn{লোকাল}} — \emph{lokāl} — \textquotedblleft{}a local train; a stopping train\textquotedblright{}.

\begin{quote}
\bn{লোকাল} \bn{ট্রেন} \bn{সব} \bn{স্টেশনে} \bn{থামে।}
\end{quote}

\begin{quote}
\emph{lokāl ṭren sôb sṭeshône thāme.}
\end{quote}

\begin{quote}
The local train stops at every station.
\end{quote}

On its own \bn{লোকাল} is the train: \bn{লোকালে} \bn{এলাম}, I came on the local.

\begin{cousinweb}
\textbf{\bn{লোকাল}} (\emph{lokāl}) is borrowed from English \emph{local}, short for \emph{local train}, the slow suburban train that stops everywhere.
\end{cousinweb}

\begin{grammarlens}[title={What you've built}]
One more word for getting around.
\end{grammarlens}

\subsection*{Your turn}
\begin{itemize}
\raggedright
  \item \emph{Say it:} \emph{lokāl}
  \item \emph{Say it:} \emph{lokāl}, once more
  \item \emph{Say it:} \emph{lokāl ṭren sôb sṭeshône thāme.} — The local train stops at every station.
  \item \emph{Recall:} say the Bengali for a compartment, then the Bengali for a berth, then say \emph{lokāl} again
\end{itemize}

\begin{tcolorbox}[breakable,colback=teal!4,colframe=teal!35!black,title={Before you move on}]
What does \bn{লোকাল} mean? (\textquotedblleft{}A local train; a stopping train\textquotedblright{}.) Say it once more.
\end{tcolorbox}
\section[\texorpdfstring{\bn{হুইসল} \bn{দেওয়া} (huisôl deoyā)}{হুইসল দেওয়া (huisôl deoyā)}]{\bn{হুইসল} \bn{দেওয়া} (huisôl deoyā) — to blow a whistle (of a train or a guard)}

\label{lesson:BN-C384-huisol-deoya}



\begin{quote}
Before the new one: say the Bengali for a berth, then the Bengali for a local train.
\end{quote}

\subsection*{You'll want to know: \bn{হুইসল} \bn{দেওয়া}}
\textbf{\bn{হুইসল} \bn{দেওয়া}} — \emph{huisôl deoyā} — \textquotedblleft{}to blow a whistle (of a train or a guard)\textquotedblright{}.

\begin{quote}
\bn{ট্রেনটা} \bn{হুইসল} \bn{দিয়ে} \bn{স্টেশন} \bn{ছাড়ল।}
\end{quote}

\begin{quote}
\emph{ṭrenṭā huisôl diye sṭeshôn chhāṛlô.}
\end{quote}

\begin{quote}
The train blew its whistle and left the station.
\end{quote}

Past (he/she …): \bn{সেই} \bn{ট্রেন} \bn{হুইসল} \bn{দিল}, \emph{sei ṭren huisôl dilô}.

Used of the engine's whistle and of the guard's whistle on the platform.

\begin{cousinweb}
\textbf{\bn{হুইসল} \bn{দেওয়া}} (\emph{huisôl deoyā}) joins \emph{whistle}, borrowed from English, to \bn{দেওয়া}, to give: you give a whistle in Bengali as you give a push.
\end{cousinweb}

\begin{grammarlens}[title={What you've built}]
One more word for getting around.
\end{grammarlens}

\subsection*{Your turn}
\begin{itemize}
\raggedright
  \item \emph{Say it:} \emph{huisôl deoyā}
  \item \emph{Say it:} \emph{huisôl deoyā}, once more
  \item \emph{Say it:} \emph{ṭrenṭā huisôl diye sṭeshôn chhāṛlô.} — The train blew its whistle and left the station.
  \item \emph{Recall:} say the Bengali for a berth, then the Bengali for a local train, then say \emph{huisôl deoyā} again
\end{itemize}

\begin{tcolorbox}[breakable,colback=teal!4,colframe=teal!35!black,title={Before you move on}]
What does \bn{হুইসল} \bn{দেওয়া} mean? (\textquotedblleft{}To blow a whistle (of a train or a guard)\textquotedblright{}.) Say it once more.
\end{tcolorbox}
